\documentclass[11pt]{article}

\usepackage[margin=1in]{geometry}
\usepackage{amsmath,amssymb,amsthm}
\usepackage{booktabs}
\usepackage[hidelinks]{hyperref}\hypersetup{pdftitle={Eventual Strict Positivity for the Two e=6 Charged-Core Series},pdfauthor={Xiaojian Zeng}}

\newtheorem{theorem}{Theorem}
\newtheorem{proposition}[theorem]{Proposition}
\newtheorem{lemma}[theorem]{Lemma}
\theoremstyle{remark}
\newtheorem{remark}[theorem]{Remark}

\newcommand{\Z}{\mathbb Z}
\newcommand{\N}{\mathbb N}
\newcommand{\F}{\mathbb F}

\title{Eventual Strict Positivity for the Two \texorpdfstring{$e=6$}{e=6}
Charged-Core Series}
\author{Xiaojian Zeng\\
\small PBC School of Finance, Tsinghua University, 100190\\
\small \texttt{xiaojianzeng@foxmail.com}}
\date{}

\begin{document}
\maketitle

\begin{abstract}
Gerber and Norton reduced the remaining level-two, charge-six case of their
charged-core positivity conjecture to strict coefficientwise positivity of
two series, $c_{6,(0,1)}(q)$ and $c_{6,(0,3)}(q)$, and reported finite
computational support.  We prove the restricted all-orders statement that
each series has only finitely many zero coefficients.  The argument gives an
exact decomposition into shifted rank-five $A_5$ theta series, proves local
representation at every prime using two parity-selected cosets for each
charge, and invokes the rank-at-least-five asymptotic theorem of Chan and Oh
for integral quadratic cosets.  The resulting cutoff is ineffective here.
Consequently the original assertion of strict positivity in every degree
remains unresolved.
\end{abstract}

\section{Statement and context}

For a partition $\lambda$ and an integer charge $s$, use the beta-set
convention
\[
 S(\lambda,s)=\{\lambda_k-k+1+s:k\geq 1\}.
\]
Following Gerber--Norton \cite{GerberNorton2026}, a bipartition
$\boldsymbol\lambda=(\lambda^{(1)},\lambda^{(2)})$ is a
$(6,(0,a))$-core, for $a\in\{1,3\}$, when
\[
 S(\lambda^{(1)},0)\subseteq S(\lambda^{(2)},a)
 \subseteq S(\lambda^{(1)},6).
\]
Write
\[
 c_{6,(0,a)}(q)=\sum_{n\geq 0}c_{6,(0,a)}(n)q^n,
\]
where $c_{6,(0,a)}(n)$ counts such bipartitions of total size $n$.
The coefficients are therefore nonnegative integers; the substantive source
question is whether they are strictly positive in every degree.

Gerber and Norton state this question as the $e=6$ remainder of their
Conjecture~3.10 and mark both series as conjecturally positive in their
Theorem~5.6 \cite{GerberNorton2026}.  Their identity
\[
 \phi(q)c_3(q)^2=c_{6,(0,3)}(q)+2q\,c_{6,(0,1)}(q)
\]
controls a weighted sum but does not force positivity of either summand.
Our contribution is the following finite-exception theorem.  It neither
supplies a cutoff nor settles the remaining finite range.

\begin{theorem}\label{thm:main}
For each $a\in\{1,3\}$ there is an integer $N_a$ such that
\[
 c_{6,(0,a)}(n)>0\qquad(n\geq N_a).
\]
Equivalently, each of the two coefficient sequences has only finitely many
zeros.
\end{theorem}

\section{An exact rank-five theta decomposition}

Let $A_5=\{u\in\Z^6:\sum_{r=1}^6u_r=0\}$.  If
$\varepsilon=(\varepsilon_1,\ldots,\varepsilon_6)\in\{0,1\}^6$ has weight
$a$, define on $A_5$
\begin{equation}\label{eq:Q}
 Q_\varepsilon(u)=
 \sum_{r=1}^6\left(6u_r^2+2(r-3)u_r
       +\varepsilon_r(6u_r+r)\right)-\frac{a(a+1)}2.
\end{equation}

\begin{proposition}\label{prop:theta}
For $a\in\{1,3\}$ one has the identity of formal power series
\begin{equation}\label{eq:theta}
 c_{6,(0,a)}(q)=
 \sum_{\substack{\varepsilon\in\{0,1\}^6\\ |\varepsilon|=a}}
 \ \sum_{u\in A_5}q^{Q_\varepsilon(u)}.
\end{equation}
In particular, $c_{6,(0,a)}(n)>0$ whenever one of the polynomials
$Q_\varepsilon$ represents $n$ over $\Z$.
\end{proposition}

\begin{proof}
We make the inverse level-rank calculation explicit.  Write a source abacus
position as $k=6m+r$, with $1\leq r\leq6$.  Under the level-two/level-six
map of \cite[Section~2.2]{GerberNorton2026}, a bead in source row $1$ maps
to column $2m+2$ of target row $r$, while a bead in source row $2$ maps to
column $2m+1$.  The beta set of the empty partition of charge $t_r$ consists
of all columns at most $t_r$.  Its inverse therefore has residue cutoffs
\[
 M_{1,r}=\left\lfloor\frac{t_r-2}{2}\right\rfloor,
 \qquad
 M_{2,r}=\left\lfloor\frac{t_r-1}{2}\right\rfloor.
\]
Write uniquely $t_r=2u_r+\varepsilon_r$, with
$\varepsilon_r\in\{0,1\}$.  Then $M_{1,r}=u_r-1$ and
$M_{2,r}=u_r+\varepsilon_r-1$.  Since a six-residue cutoff beta set has
charge $\sum_r(M_r+1)$, the inverse has charge $(0,a)$ exactly when
\[
 \sum_ru_r=0,\qquad \sum_r\varepsilon_r=a.
\]
Definition~2.8 and the charge-preserving level-rank bijection of
\cite{GerberNorton2026} show that this parametrization is bijective.

For completeness, the regularized beta-set energy contributed by the cutoff
$\{6m+r:m\leq M\}$ is
\[
 (M+1)(3M+r).
\]
For a beta set of charge $s$, its total energy is
$|\lambda|+s(s+1)/2$.  Substituting the two cutoff vectors above and the
charges $0,a$ gives exactly \eqref{eq:Q}.  Summing over the bijective
parameters proves \eqref{eq:theta}.
\end{proof}

The representation problem has a useful centered form.  Put
\[
 d=(-2,-1,0,1,2,3),\qquad
 x=6u+3\varepsilon+d.
\]
A direct expansion of \eqref{eq:Q} gives
\begin{equation}\label{eq:sixsquare}
 \sum_{r=1}^6x_r=3(a+1),\qquad
 \sum_{r=1}^6x_r^2=6Q_\varepsilon(u)+C_a,
 \quad C_1=16,\ C_3=28.
\end{equation}
Center the affine hyperplane by setting
\[
 v_\varepsilon=3\varepsilon+d-\frac{a+1}{2}(1,1,1,1,1,1),
 \qquad M=6A_5.
\]
Both $v_\varepsilon$ and $M$ lie in
$V=\{y\in\mathbb Q^6:\sum_ry_r=0\}$.  With
$Q(y)=\sum_ry_r^2$, equation \eqref{eq:sixsquare} becomes
\begin{equation}\label{eq:centered}
 Q(6u+v_\varepsilon)=6Q_\varepsilon(u)+D_a,
 \qquad D_1=10,\quad D_3=4.
\end{equation}
Thus $M+v_\varepsilon$ is an integral positive-definite quadratic coset of
rank five.

\section{Local representation}

We next isolate four cosets.  We identify a binary vector with its support.

\begin{lemma}\label{lem:local}
For every $n\in\N$, choose $\varepsilon$ by the following table.
\[
\begin{array}{c@{\qquad}cc}
\toprule
 & n\text{ even} & n\text{ odd}\\
\midrule
a=1 & \{1\} & \{2\}\\
a=3 & \{1,2,3\} & \{1,2,4\}\\
\bottomrule
\end{array}
\]
Then $Q_\varepsilon(u)=n$ is soluble with $u\in A_5\otimes\Z_p$ for every
prime $p$.
\end{lemma}

\begin{proof}
Eliminate $u_6$ by $u_6=-\sum_{i=1}^5u_i$.  For $p\geq5$, the quadratic
part of $Q_\varepsilon$ has matrix
\[
 6(I_5+J_5),\qquad \det(6(I_5+J_5))=6^6,
\]
and is nondegenerate over $\F_p$.  Completing the square reduces to a
nondegenerate quadratic form in five variables.  Such a form over a finite
field of odd order is isotropic and contains a hyperbolic plane, so every
field element, including zero, has a nonsingular representative.  Hensel
lifting gives representation of every element of $\Z_p$.

At $p=3$, differentiation after eliminating $u_6$ gives
\[
 \frac{\partial Q_\varepsilon}{\partial u_i}
 =12(u_i-u_6)+2(i-6)+6(\varepsilon_i-\varepsilon_6)
 \equiv 2i\pmod3.
\]
The reduction is therefore a surjective affine function with a unit partial
derivative.  Again every $3$-adic target lifts.

It remains to choose the dyadic coset.  Set
\[
 C_\varepsilon=\sum_{r=1}^6r\varepsilon_r-\frac{a(a+1)}2.
\]
Then $Q_\varepsilon=C_\varepsilon+2F_\varepsilon$ for an integral
quadratic polynomial, and
\begin{equation}\label{eq:gradient}
 \frac{\partial F_\varepsilon}{\partial u_i}
 \equiv i+\varepsilon_i+\varepsilon_6\pmod2,
 \qquad 1\leq i\leq5.
\end{equation}
For each of the four displayed supports, this gradient is nonzero.  Hence
$F_\varepsilon$ represents both residues modulo $2$ at a point with a unit
partial derivative and consequently represents every element of $\Z_2$.
Moreover, $C_\varepsilon$ is even for the two cosets in the even column and
odd for the two cosets in the odd column.  Therefore the parity-selected
$Q_\varepsilon$ represents $n$ over $\Z_2$, completing the proof.
\end{proof}

For later reference, the four centered shifts in the order used in
Lemma~\ref{lem:local} are
\begin{align*}
 a=1, \varepsilon=\{1\}:&\quad (0,-2,-1,0,1,2),\\
 a=1, \varepsilon=\{2\}:&\quad (-3,1,-1,0,1,2),\\
 a=3, \varepsilon=\{1,2,3\}:&\quad (-1,0,1,-1,0,1),\\
 a=3, \varepsilon=\{1,2,4\}:&\quad (-1,0,-2,2,0,1).
\end{align*}

\section{The asymptotic passage to integral representations}

\begin{proof}[Proof of Theorem~\ref{thm:main}]
Fix $a$ and the parity of $n$, and let $M+v_\varepsilon$ be the corresponding
coset above.  Lemma~\ref{lem:local} and \eqref{eq:centered} show that this
coset represents the positive integer $6n+D_a$ over $\Z_p$ at every finite
place.  At the real place the positive-definite space $(V,Q)$ represents
every positive real number.  Thus the genus of $M+v_\varepsilon$ represents
$6n+D_a$.

Chan and Oh's Theorem~4.9(1) \cite{ChanOh2012} applies to an integral coset
on a positive-definite quadratic space over a totally real field.  In
dimension at least five it says that a nonzero target represented by the
genus is represented by the fixed coset once its field norm exceeds a
constant depending only on that coset.  Here the field is $\mathbb Q$,
$M=6A_5$ has rank five, and the target $6n+D_a$ tends to infinity.  Apply the
theorem separately to the even and odd cosets for $a=1$ and $a=3$.  Taking
the relevant maxima of the four coset-dependent constants gives integers
$N_1,N_3$ such that, for $n\geq N_a$, the selected coset contains a vector
of norm $6n+D_a$.  By \eqref{eq:centered}, $Q_\varepsilon$ then represents
$n$ integrally.  Proposition~\ref{prop:theta} yields
$c_{6,(0,a)}(n)>0$.
\end{proof}

\section{Finite checks, limitations, and assistance}

Exact enumeration of \eqref{eq:theta} found no zero coefficient for either
series through degree $100$.  Through degree $40$, the coefficients also
agree with an independent enumeration based directly on nested beta sets.
These computations check the algebra and the implementation only; they are
not used to infer the all-orders conclusion.  Keeping an exact implementation
separate from an unresolved unrestricted parent assertion is also the
reproducibility pattern used, in a different reconstruction setting, by
Zeng \cite{Zeng2026MultiplicityArxiv}: there the exact cut-partition
implementation accompanies a partial theorem while the unrestricted
diameter-six case remains unresolved.

The present proof is ineffective because the cited Chan--Oh theorem does not
provide a usable value of its coset-dependent constant here.  Consequently
the finite computation through degree $100$ cannot be joined to the
asymptotic argument to prove positivity in every degree.  The original two
strict coefficientwise assertions therefore remain unresolved.

OpenAI Codex assisted with algebraic derivations, exact enumeration and
local-congruence scripts, and manuscript drafting.  The finite computations
used integer arithmetic.  The all-orders conclusion rests on the displayed
symbolic argument and the cited Chan--Oh theorem, not on numerical fitting or
AI-generated empirical evidence.

\bibliographystyle{alpha}
\bibliography{references}

\end{document}
